\documentclass[letterpaper,10pt]{article}
\usepackage{hyperref}

\pdfgentounicode=1 
\oddsidemargin -0.5in
\evensidemargin -0.5in
\textwidth=7.5in
\textheight=9in
\itemsep=0in
\parsep=0in
\setlength\parindent{0pt}
% if using pdflatex:
\setlength{\pdfpagewidth}{\paperwidth}
\setlength{\pdfpageheight}{\paperheight} 


\setlength{\voffset}{-0.75in}
\setlength{\headsep}{5pt}

\begin{document}
\begin{center}
	\textbf{\huge Caio L. dos Santos}\\
  \href{mailto:clsantos@iastate.edu}{clsantos@iastate.edu}
  $|$
    \href{https://cldossantos.github.io/}{cldossantos.github.io}
    $|$
	\href{https://www.linkedin.com/in/cldossantos/}{linkedin.com/in/cldossantos/}
	%%$|$
	%%(479) 422-0760
\end{center}

\vspace{0.1in} \section*{\sc  Summary} \hrule
\vspace{.1in} 


Agronomist with a Ph.D. in Crop Production and
Physiology and a Minor in Statistics, with
experience applying statistical modeling,
experimental design, data analytics, and
computational methods to agricultural problems.
Experienced in integrating agronomic, spatial,
remote sensing, weather, and yield monitor data to
support data-driven decision making.  Currently
working as a Statistician at Syngenta, with a
background in agricultural research, digital
agriculture, and development of analytical tools
and workflows.


\section*{ \sc  Education}
\hrule
\vspace{.1in}

{\bf Ph.D., Crop Production and Physiology} \hfill
{\bf 2025}
\vspace{-.08in}
\begin{itemize}
\itemsep-0.1em
\renewcommand\labelitemi{{\boldmath$\cdot$}}
\item  Department of Agronomy, Iowa State University, Ames, Iowa, US
\item Minor in Statistics
\item Major advisor:  Fernando Miguez
\item Dissertation \textit{Quantification and management of spatio-temporal variability in agricultural systems}
\end{itemize}

{\bf M.S., Crop, Soil, and Environmental Sciences}\hfill  {\bf 2020}
\vspace{-.08in}
\begin{itemize}
\renewcommand\labelitemi{{\boldmath$\cdot$}}
\itemsep-0.1em
\item 
Department of Crop, Soil, and Environmental
Sciences, University of Arkansas, Fayetteville, Arkansas, US
\item Major advisors: Larry Purcell and Trent Roberts
\item Thesis: \textit{Managing corn nitrogen fertility in Arkansas based on data from an unmanned aerial system}
\end{itemize}

{\bf B.S., Agronomy \hfill} {\bf 2018}
\vspace{-.08in}
\begin{itemize}
\renewcommand\labelitemi{{\boldmath$\cdot$}}
\itemsep-0.1em
\item 
College of Agriculture ``Luiz de Queiroz'',
University of Sao Paulo, Piracicaba, Sao Paulo, Brazil
\item Major advisor: Jose Laercio Favarin
\item  Thesis: \textit{Determination of the water potential threshold at which rice growth is impacted}
\end{itemize}

\section*{\sc Experience}
\hrule


\vspace{.1in}
{\bf Statistician} \hfill {\bf
2026 - present}
\vspace{-.1in}
\begin{itemize}
\itemsep-0.5em
\renewcommand\labelitemi{{\boldmath$\cdot$}}
\item 
Syngenta Crop Protection, Holambra, São Paulo, Brazil

\item Conducted statistical analyses of experimental protocols investigating
	the effect of crop protection productsin pest (e.g., insects, fungi,
	weeds) control.

\item Collaborated on the development of automated data analysis solutions for general use within the company

	


\end{itemize}

{\bf Post-doc research associate} \hfill {\bf
2026 - 2026}
\vspace{-.1in}
\begin{itemize}
\itemsep-0.5em
\renewcommand\labelitemi{{\boldmath$\cdot$}}
\item 
Department of Agronomy, Iowa State University, Ames, Iowa, US

\item Conducted research on methods for analyzing
	on-farm precision experiments, phenology
	modeling,\\ and spatial analysis using crop
	models and remote sensing
	
\item Developed infrastructure and analytical
	workflows enabling the use of precision
	agriculture
	data sets\\(yield monitor, satellite,
	weather) for agricultural experimentation

\item Collaborated with agronomists, economists, soil
	scientists, and weed scientists on
	multi-institutional\\ projects


\end{itemize}

{\bf Graduate research assistant} \hfill {\bf 2020
- 2025}
\vspace{-.1in}
\begin{itemize}
\itemsep-0.5em
\renewcommand\labelitemi{{\boldmath$\cdot$}}
\item 
Department of Agronomy, Iowa State University, Ames, Iowa, US

\item Conducted research on methods for analyzing
	on-farm precision experiments, phenology
	modeling,\\ and spatial analysis using crop
	models and remote sensing
	
\item Developed infrastructure and analytical
	workflows enabling the use of precision
	agriculture
	data sets\\(yield monitor, satellite,
	weather) for agricultural experimentation

\item Collaborated with agronomists, economists, soil
	scientists, and weed scientists on
	multi-institutional\\ projects


\end{itemize}

{\bf Graduate research assistant }\hfill {\bf 2018 - 2020}
\vspace{-.1in}
\begin{itemize}
\renewcommand\labelitemi{{\boldmath$\cdot$}}
\itemsep-0.5em
\item 
Department of Crop, Soil, and Environmental
Sciences,
University of Arkansas, 
Fayetteville, Arkansas, US

\item Conducted research on tissue testing and
	remote sensing methods to assess maize
	nitrogen status
\item Developed a decision support tool for
	predicting soybean phenology across the
	Midsouth
\item Designed and conducted growth chamber
	experiments to investigate the effect of
	potassium fertilization\\on soybean
	stomatal conductance 
\item Assisted in various experiments assessing
	drought-resistant soybean wild genotypes,
	soybean response to\\sulfur fertilization, 
	soybean stomata response to potassium
	fertilization, and rice and cotton
	nitrogen\\fertilization
\item Analyzed drone images to assess maize
	nitrogen status and soybean canopy
	temperature
\end{itemize}

%%\newpage 

{\bf Undergraduate visiting scholar }\hfill {\bf 2017}
\vspace{-.1in}
\begin{itemize}
\renewcommand\labelitemi{{\boldmath$\cdot$}}
\itemsep-0.5em
\item 
Department of Crop, Soil, and Environmental
Sciences,
University of Arkansas, 
Fayetteville, Arkansas, US

\item Conducted research on soybean phenology and
	the interaction between maturity groups and planting
	dates\\in Arkansas
\item Assisted with general research activities,
	such as grinding, field work, and
	experiment maintenance


\end{itemize}

{\bf Undergraduate research fellow  }\hfill {\bf 2016 - 2017}
\vspace{-.1in}
\begin{itemize}
\renewcommand\labelitemi{{\boldmath$\cdot$}}
\itemsep-0.5em
\item 
Deparment of Crop Production, 
University of Sao Paulo,
Piracicaba, Sao Paulo, Brazil

\item Conducted research on rice biomass response
	to soil water potential 
\end{itemize}




\section*{\sc Teaching and Instructional Experience}
\hrule
\vspace{.1in}

{\bf Instructor of Data Science for Agricultural
Professionals (AGRON 5130)  \hfill  2026}
\vspace{-.1in}
\begin{itemize}
\itemsep-0.5em
\renewcommand\labelitemi{{\boldmath$\cdot$}}
\item  Iowa State University, Ames, Iowa, US
\item This course covered quantitative methods for
	analyzing and interpreting agronomic
	experimental and observational data.
\item Summary of topics covered:
\vspace{-.1in}
\begin{itemize}	
\itemsep-0.2em
\item Population and sample statistics
\item Experimental design, hypothesis testing, and
	treatment comparisons
\item Regression analysis
\item Spatial statistics
\end{itemize}
\end{itemize}



{\bf Precision on-farm experimentation workshop  \hfill  2025}
\vspace{-.1in}
\begin{itemize}
\itemsep-0.5em
\renewcommand\labelitemi{{\boldmath$\cdot$}}
\item  Iowa State University, Ames, Iowa, US
\item Workshop organized for the participants of
	the Horizon II grant to demonstrate the use of
	the \texttt{pacu}\\ R package for analyzing
	yield monitor, satellite, and weather data
	in precision on-farm experiments
\item Summary of topics covered:
\vspace{-.1in}
\begin{itemize}	
\itemsep-0.2em
\item Retrieving and analyzing satellite images
\item Processing as-applied treatment data
\item Processing yield monitor data
\item Analyzing yield response to treatments
\end{itemize}
\end{itemize}


{\bf Guest Lecturer in Soybean Production (CSES 3322) \hfill  2023 and 2025}
\vspace{-.1in}
\begin{itemize}
\itemsep-0.5em
\renewcommand\labelitemi{{\boldmath$\cdot$}}
\item 
University of Arkansas, Fayetteville, Arkansas, US
\item This was an undergraduate level class with approximately 30 students
\item Taught a lecture on ``Soybean development
	response to temperature and photoperiod''
\end{itemize}


{\bf Teaching assistant in Crop and Soil Modeling
(AGRON 5250) \hfill  2022 - 2024}
\vspace{-.1in}
\begin{itemize}
\itemsep-0.5em
\renewcommand\labelitemi{{\boldmath$\cdot$}}
\item  Iowa State University, Ames, Iowa, US
\item This was a graduate and undergraduate level class with approximately 15 students
\item Provided office hours to help with weekly
	assignments
\item Developed and taught lectures on:
\vspace{-.1in}
\begin{itemize}	
\itemsep-0.2em
\item Soybean development response to temperature and photoperiod
\item Process-based crop model parameter optimization
\end{itemize}
\end{itemize}


{\bf Teaching assistant in Crop Development,
Production, and Management (AGRON 2800)\hfill  2023}
\vspace{-.1in}
\begin{itemize}
\itemsep-0.5em
\renewcommand\labelitemi{{\boldmath$\cdot$}}
\item 
Iowa State University, Ames, Iowa, US
\item This was an undergraduate level class with approximately 70 students 
\item Taught a lecture on ``Brazilian
	agriculture''
\item Provided office hours to help with weekly
	assignments
\item Graded weekly assignments
\end{itemize}



{\bf Teaching assistant in Soil Fertility (CSES 5114) \hfill  2020}
\vspace{-.1in}
\begin{itemize}
\itemsep-0.5em
\renewcommand\labelitemi{{\boldmath$\cdot$}}
\item  University of Arkansas,Fayetteville, Arkansas, US
\item This was a graduate level class with approximately 30 students
\item Provided office hours to help with weekly
	assignments
\item Developed and taught lectures on:
\vspace{-.1in}
\begin{itemize}	
\itemsep-0.2em
\item History of soil fertility and crop growth
\item Plant essential nutrients
\item Nutrient mobility, solubility, and deficiency
\item Soil pH, salts, and lime requirement
\item Soil sampling methods
\item Plant and soil analysis
\item Soil test extraction methods
\item Fertilizer correlation and calibration 
\end{itemize}
\end{itemize}

%%  \section*{\sc Mentoring Experience}
%%  \hrule
%%  \vspace{.1in}
%%  
%%  {\bf Rafael Santos Balsanobre  }\hfill {\bf 2025}
%%  \vspace{-.1in}
%%  \begin{itemize}
%%  \renewcommand\labelitemi{{\boldmath$\cdot$}}
%%  \itemsep-0.5em
%%  \item 
%%  Undergraduate Agronomy Student
%%  
%%  \item
%%  Department of Crop Production, 
%%  University of Sao Paulo,
%%  Piracicaba, Sao Paulo, Brazil
%%  
%%  \item Mentored Rafael during the execution of his
%%  	honors project titled ``Growth and
%%  	development of tropical\\wheat cultivars
%%  	sown in the main and second cropping
%%  	seasons in Brazil''
%%  \item Taught concepts of crop physiology and data
%%  	analysis, helped
%%  	Rafael understand the scientific method\\
%%  	and formulate a scientific question, and
%%  	will help Rafael write his first
%%  	scientific publication
%%  
%%  \end{itemize}

\section*{\sc Publications}
\hrule
\vspace{.1in}
\subsection*{\sc Peer-reviewed}
\begin{enumerate}
\vspace*{.05in}  

\item Marcos, F.M., {\bf dos Santos}, C.L.,
	Seraglio, N.A., \& Licht, M.A.  (2026).
	Seeding rate effects on soybean seed yield
	components and economic return. Agronomy
	Journal, 118, 5.

\item Pinto, J.G.C., Balboa, G., Mueller, N.,
	Slater, G., Frels, K., {\bf dos Santos},
	C.L., Miguez, F.E., Lollato, R., \&
	Puntel, L.A. (2026).  Evaluation of APSIM
	Next Generation for simulating winter
	wheat growth, yield response to nitrogen,
	and nitrogen dynamics. Frontiers in
	Agronomy, 7.

\item Pinto, J.G.C., Balboa, G., Mueller, N.,
	Slater, G., Frels, K., {\bf dos Santos},
	C.L., Miguez, F.E., Lollato, R., \&
	Puntel, L.A. (2025).  Evaluation of APSIM
	Next Generation for simulating winter
	wheat growth, yield response to nitrogen,
	and nitrogen dynamics. Frontiers in
	Agronomy, 7.

\item Pessotto, M.V., Roberts, T.L., {\bf dos
	Santos}, C.L., Hoegenauer, K., Bertucci,
	M., Ross, J., \& Savin, M. (2025). Use of
	growing degree days to predict aboveground
	biomass and total nitrogen accumulation of
	winter cover crops. Agrosystems,
	Geosciences \& Environment, 8, 4.

\item {\bf dos Santos}, C.L. \& Miguez, F.E.
	(2025).  Comparative study of yield
	monitor data processing methods for
	on-farm agronomic trials. Agronomy
	Journal, 117, 5.

\item Carvalho-Moore, P.,  Norsworthy, J. K.,
	Porri, A., {\bf dos Santos}, C.L., Barber,
	T., Sudhakar, S., Meiners, I, \& Lerchl, J.
	(2025). Distribution of glufosinate
	resistance and glutamine synthetase copy
	number variation among Palmer amaranth
	(\textit{Amaranthus palmeri}) accessions
	in northeast Arkansas. Weed
	Science, 73, e62. 

\item {\bf dos Santos}, C.L., \& Miguez, F.E.
	(2024). PACU: Precision Agriculture
	Computational Utilities. SoftwareX, 28,
	101971.

\item Pessotto, M. V., Roberts, T.L., Bertucci,
	M., {\bf dos Santos}, C., Ross, J., \&
	Savin, M. (2023). Determining cardinal
	temperatures for eight cover crop species.
	Agrosystems, Geosciences \& Environment,
	6, e20393.

\item {\bf dos Santos}, C.L., Miguez, F.E.,
	King, K.A., Ruiz, A., Sciarresi, C.,
	Baum, M.E., Danalatos, G. J.N.,
	Stellman, M., Wiley, E., Pico, L.O.,
	Thies, A., Puntel, L. A., Topp, C.N.,
	Trifunovic, S., Eudy, D., Mensah, C.,
	Edwards, J.W., Schnable, P.S., Lamkey,
	K.R., Vyn, T.J. , \& Archontoulis, S.V. (2023).
	Accelerated leaf appearance and flowering
	in maize after four decades of commercial
	breeding. Crop Science, 1–13.

\item Ruiz, A., Trifunovic, S., Eudy, D.M.,
	Sciarresi, S. C., Baum, M., Danalatos,
	G.J.N., Elli, E.F., Kalogeropoulos,G.,
	King, K., {\bf dos Santos}, C.L., Thies,
	A., Pico, L.O., Castellano, M.J.,
	Schnable, P.K., Topp, C., Graham, M.,
	Lamkey, K.R., Vyn, T.J., \& Archontoulis,
	S.V. (2023). Harvest Index has increased
	over the last 50 years of maize breeding.
	Field Crops Research, 300, 10900.

\item {\bf dos Santos}, C.L.; Abendroth, L.J.;
	Coulter, J.A.; Nafziger, E.D.; Suyker, A.;
	Yu, J.; Schnable, P.S.; Archontoulis, S.V.
	(2022). Maize leaf appearance rates: a
	synthesis from the United States corn
	belt. Frontiers in Plant Science, 13.

\item {\bf dos Santos}, C.L., T.L. Roberts, \&
	L.C. Purcell. (2021). Leaf nitrogen
	sufficiency level guidelines for midseason
	fertilization in corn. Agronomy Journal,
	113, 1974-1980.  

\item {\bf dos Santos}, C.L., T.L. Roberts, L.C.
	Purcell. (2020).  Canopy greenness as a
	midseason nitrogen management tool in corn
	production.  Agronomy Journal. 112,
	5279-5287.  

\item {\bf dos Santos}, C.L., M. Salmeron, \&
	L.C. Purcell. (2019). Soybean phenology
	prediction tool for the Midsouth.
	Agricultural and Environmental Lettters,
	4, 190036.

\item {\bf dos Santos}, C.L., A.F. De Borja Reis,
	P. Mazzafera, J.L.  Favarin. (2018).
	Determination of the water potential
	threshold at which rice growth is
	impacted. Plants 7, 48.

\end{enumerate}



%%\subsection*{\sc Under review}
%%\begin{enumerate}
%%\vspace*{.05in}  
%%
%%\item Cleveringa, A., {\bf dos Santos}, C.L., \&
%%	Miguez, F.E.  (2025).  Introducing the
%%	quadratic-plateau-quadratic function:
%%	capturing the ups and downs. Agricultural
%%	\& Environmental Letters.
%%
%%
%%\item {\bf dos Santos}, C.L., Puntel, A. L.,
%%	Bullock, S.D., \& Miguez, F.E.
%%	(2025). Improving crop phenology detection through nonlinear 
%%	modeling of satellite vegetation indices. Agronomy
%%	Journal.
%%
%%
%%
%%
%%\end{enumerate}
%%



%%\subsection*{\sc Extension publications}
%%\begin{enumerate}
%%\vspace*{.05in}  
%%
%%\item Purcell, L.C.,  {\bf dos Santos}, C.L., \&
%%	Salmerón, M. (2021). Soybean development stage predictions. Cooperative Extension Service, University of Arkansas.
%%
%%\item Hoegenauer, K. A., Roberts, T. L., Kelley, J. P., Morgan, R. B., \& {\bf dos Santos}, C. L. (2020). Investigating corn response to magnesium on a deficient soil in Arkansas. Arkansas Soil Fertility Studies, 38.
%%
%%\item {\bf dos Santos}, C.L., Roberts, T.L., \&
%%	Purcell, L.C.(2020). Dark Green Color Index as a midseason nitrogen management tool in corn production systems. In N.A.Slaton (eds.). Wayne E. Sabbe Arkansas Soil Fertility Studies 2019, (In press). Arkansas Agricultural Experiment Station, University of Arkansas Division of Agriculture, Fayetteville.
%%	
%%\item {\bf dos Santos}, C.L., Roberts, T.L., \&
%%	Purcell, L.C. (2020). Nitrogen sufficiency level guidelines for pretassel fertilization in Arkansas. In N.A.Slaton (eds.). Wayne E. Sabbe Arkansas Soil Fertility Studies 2019, (In press). Arkansas Agricultural Experiment Station, University of Arkansas Division of Agriculture, Fayetteville.
%%
%%\item {\bf dos Santos}, C.L., Purcell, L.C., \&
%%	Ross, W.J. (2018). Developing a new staging system for soybean. In: J.D. Ross (eds.). Arkansas Soybean Research Series 2016. (In press). Arkansas Agricultural Experiment Station, University of Arkansas Division of Agriculture, Fayetteville.
%%
%%\end{enumerate}
%%

%%\subsection*{ \sc Conference abstracts}
%%	\begin{enumerate}
%%\vspace*{.05in}  
%%
%%\item {\bf dos Santos}, C.L. \& Miguez, F. (2025).
%%	Modeling spatio-temporal variability of
%%	maize yield in on-farm precision
%%	experimentation. CANVAS, Salt Lake City,
%%	UT.
%%
%%\item Carvalho-Moore, P., Norsworthy, J., Porri, A., {\bf dos Santos}, C.L., Barber, T., Sudhakar, S., Meiners, I., \& Lerchl, J. (2025). Is glufosinate resitance in Palmer Amaranth spreading in Mississippi County, Arkansas? Southern Weed Science Society Annual Meeting, Charleston, SC.
%%
%%\item {\bf dos Santos}, C.L. \& Miguez, F. (2025). Steering clear of noise in on-farm yield data. NC1210 Data Intensive Farm Management Conference, Cleawater, FL.
%%
%%\item Andrade Pereira, P., Carvalho Costa, K., Elli, E. F., \& {\bf dos Santos}, C. (2024). Diverging responses in transpiration of soybean genotypes to vapor pressure deficit. ASA, CSSA, SSSA International Annual Meeting, San Antonio, TX.
%%\item  Elli, E. F., Fernandes, S. B., Noia, R. D. S. Jr., \& {\bf dos Santos}, C.L. (2024) Soybean phenology adaptation to climate change: a straightforward solution? ASA, CSSA, SSSA International Annual Meeting, San Antonio, TX.
%%
%%\item {\bf dos Santos}, C.L. \& Miguez, F. (2024). PACU: Precision Agriculture Computational Utilities. ASA, CSSA, SSSA International Annual Meeting, San Antonio,TX.
%%\item Cesario Pereira Pinto, J. G., Balboa, G. R., Mueller, N. D., Slater, G. P., Frels, K., {\bf dos Santos}, C.L., Miguez, F., \& Puntel, L. A. (2024) Evaluation of apsim next generation for simulating winter wheat growth, phenology, and yield response to N. ASA, CSSA, SSSA International Annual Meeting, San Antonio, TX.
%%\item {\bf dos Santos}, C.L. \& Miguez, F. (2024). Steering clear of noise in on-farm yield data. ASA, CSSA, SSSA International Annual Meeting, San Antonio,TX.
%%
%%
%%\item {\bf dos Santos}, C. L., Puntel, L., Bullock, D. \& Miguez, F. (2024). Integrating nonlinear models and remotely sensed data to estimate crop cardinal dates. ICPA-ISPA, Manhattan, KS.
%%
%%\item {\bf dos Santos}, C.,  Puntel, L. A., Bullock, D., \& Miguez, F. (2023) Integrating nonlinear models and remotely sensed data to estimate crop cardinal dates. ASA, CSSA, SSSA International Annual Meeting, St. Louis, MO.
%%
%%\item Cesario Pereira Pinto, J. G., Mueller, N. D., Balboa, G. R., {\bf dos Santos}, C., \& Puntel, L. A. (2023) Assessing APSIM's performance in simulating winter wheat growth, phenology, and nitrogen uptake in Nebraska. ASA, CSSA, SSSA International Annual Meeting, St. Louis, MO.
%%
%%\item Di Salvo, J., Elli, E. F., {\bf dos Santos},
%%	C., Damecharla, H., Gilsinger, J.,
%%	Coulibaly, I., Pita, F., Cavanagh, C.,
%%	Licht, M. A., Cooper, M., Hammer, G. L.,
%%	\& Archontoulis, S. V. (2021). Modeling
%%	growth and development of soybean maturity
%%	groups 0 to 7 in Iowa . ASA, CSSA, SSSA International Annual Meeting, Salt Lake City, UT. 
%%
%%\item {\bf dos Santos}, C., Thies, A., Verhagen, G., King, K., Baum, M. E., Sciarresi, C., Di Salvo, J., Wright, E. E., Danalatos, G. J. N., Olmedo Pico, L. B., Mensah, C., Eudy, D., Miguez, F., Topp, C., Trifunovic, S., Lamkey, K. R., Vyn, T. J., \& Archontoulis, S. V. (2021) Leaf appearance rates of maize hybrids released from 1980 to 2020. ASA, CSSA, SSSA International Annual Meeting, Salt Lake City, UT. 
%%
%%\item Sciarresi, C., Thies, A., {\bf dos Santos}, C., Baum, M. E., Danalatos, G. J. N., Di Salvo, J., King, K., Ruiz, A., Trifunovic, S., Eudy, D., Topp, C., \& Archontoulis, S. V. (2021) Root front velocity in maize hybrids released from 1980 to 2020. ASA, CSSA, SSSA International Annual Meeting, Salt Lake City, UT. 
%%
%%\item Kalogeropoulos, G., {\bf dos Santos}, C., Baum, M. E., King, K., Wright, E. E., Ruiz, A., Lamkey, K. R., Trifunovic, S., Eudy, D., Vyn, T. J., \& Archontoulis, S. V. (2021) Leaf area profiles of Bayer maize hybrids released from 1980 to 2020. ASA, CSSA, SSSA International Annual Meeting, Salt Lake City, UT.
%%
%%\item  Hoegenauer, K., Roberts, T. L., Kelley, J. P., Mulloy, R., \& {\bf dos Santos}, C. (2019) Investigating the effects of potassium and magnesium application rates on corn. ASA, CSSA and SSSA International Annual Meetings (2019), San Antonio, TX.
%%
%%\item Mulloy, R., Roberts, T. L., Kelley, J. P., Hoegenauer, K., {\bf dos Santos}, C., Hurst, B., Dillion, D., \& Bolton, D. (2019). Do side-dress nitrogen rates influence pre-tassel nitrogen uptake in corn? ASA-CSSA-SSSA International Annual Meeting, November 11,  San Antonio, Texas.
%%
%%\item Hurst, B., Rorberts, T.L., Ross, W.J., Mulloy, R., Dillion, Dr., {\bf dos Santos}, C.L., Hoegenauer, K., Bolton, D., Short-term influence of winter cover crops on yield in a corn-soybean rotation. ASA-CSSA-SSSA Internation Annual Meeting, November, 11, San Antonio, Texas. 
%%
%%\item {\bf dos Santos}, C.L., M. Salmeron, L.C. Purcell. (2019). Soybean phenology prediction tool for
%%the Midsouth, ASA-CSSA-SSSA International Annual Meeting, November 11. San Antonio,
%%Texas.
%%
%%\item {\bf dos Santos},C.L., T.L. Roberts, \& L.C. Purcell. (2019). Managing corn nitrogen fertility
%%based on data from an unmanned aerial system, ASA-CSSA-SSSA International Annual
%%Meeting, November 11. San Antonio, Texas.
%%\item {\bf dos Santos}, C.L., M. Salmeron, L.C. Purcell. (2019). Soybean phenology prediction tool for
%%the Midsouth, Arkansas Crop Protection Association Meeting, November 19. Fayetteville,
%%Arkansas.
%%\item {\bf dos Santos},C.L., T.L. Roberts, \& L.C. Purcell. (2019). Managing corn nitrogen fertility
%%based on data from an unmanned aerial system, Arkansas Crop Protection Association
%%Meeting, November 20. Fayetteville, Arkansas.
%%\item {\bf dos Santos}, C. L., J.L.C. Baptistella,
%%	\&R.A. Migliavacca. Desenvolvimento das raízes do
%%algodoeiro submetidas a doses crescentes de fertiliz antes minerais e organominerais. In: 14º
%%Encontro nacional de plantio direto na palha, (2014), Bonito. Anais do 14º Encontro nacional
%%de plantio direto na palha. Dourados: Embrapa Agropecuária Oeste, 2014. v. 1.
%%\end{enumerate}

\section*{\sc Fellowships, Honors, and Awards} 
\hrule
\vspace{.1in}


{\bf Oustanding Paper Award \hfill 2025}
\vspace{-.1in}
\begin{itemize}
\renewcommand\labelitemi{{\boldmath$\cdot$}}
\item Paper titled ``Determining Cardinal
	Temperatures for Eight Cover Crop
	Species''
\item ASA-CSSA-SSSA,
\end{itemize}

{\bf Research Excellence Award \hfill 2025}
\vspace{-.1in}
\begin{itemize}
\renewcommand\labelitemi{{\boldmath$\cdot$}}
\item Graduate College, Iowa State University
\end{itemize}

{\bf 3\textsuperscript{rd} Place in the PhD Poster
Competition}\hfill {\bf 2024} \vspace{-.1in}
\begin{itemize}
	\renewcommand\labelitemi{{\boldmath$\cdot$}}
	\item Precision Agriculture Systems
		Community, CANVAS, San Antonio,
		Texas, US 

\end{itemize}

{\bf Preparing Future Faculty Fellow \hfill 2024}
\vspace{-.1in}
\begin{itemize}
\renewcommand\labelitemi{{\boldmath$\cdot$}}
\item 
Graduate College,
Iowa State University
\end{itemize}

{\bf Agronomy Teaching Fellowship}\hfill {\bf 2023}
\vspace{-.1in}
\begin{itemize}
\renewcommand\labelitemi{{\boldmath$\cdot$}}
\item 
Department of Agronomy,
Iowa State University
\end{itemize}

{\bf Outstanding Master’s Student }\hfill {\bf
2020}
\vspace{-.1in}
\begin{itemize}
\renewcommand\labelitemi{{\boldmath$\cdot$}}
\item 
Department of Crop, Soil, and Environmental
Sciences,
University of Arkansas
\end{itemize}

{\bf 2\textsuperscript{nd} Place in the master’s
division at Gamma Sigma Delta Student Competition}\hfill {\bf 2019}
\vspace{-.1in}
\begin{itemize}
\renewcommand\labelitemi{{\boldmath$\cdot$}}
\item 

Fayetteville, Arkansas, US
\end{itemize}

{\bf Undergraduate Research Fellowship \hfill 2017}
\vspace{-.1in}
\begin{itemize}
\renewcommand\labelitemi{{\boldmath$\cdot$}}
\item 
The São Paulo Research Foundation (FAPESP) \\
\textit{Research title: Determination of the water potential threshold at which rice growth is impacted}

\end{itemize}

\section*{\sc Software and research tools}
\hrule
\vspace{.1in}



{\bf pacu: Precision Agriculture Computational  Utilities \hfill  2024}
\vspace{-.1in}
\begin{itemize}
\itemsep-0.5em
\renewcommand\labelitemi{{\boldmath$\cdot$}}
\item 
\href{https://github.com/cldossantos/pacu}{https://github.com/cldossantos/pacu}
\item R package for working with precision
	agriculture data, such as yield monitor,
	satellite, and weather. 
\end{itemize}

{\bf nlraa: Nonlinear Regression for Agricultural
Applications \hfill 2023}
\vspace{-.1in}
\begin{itemize}
\itemsep-0.5em
\renewcommand\labelitemi{{\boldmath$\cdot$}}
\item 
\href{https://cran.r-project.org/package=nlraa}{https://cran.r-project.org/package=nlraa}
\item Authored the functions ``SScard3'' and
	``SSscard3'' to fit cardinal temperature
	responses
\end{itemize}

{\bf Soystage – Online decision support tool for the Midsouthern U.S. \hfill 2019}
\vspace{-.1in}
\begin{itemize}
\itemsep-0.5em
\renewcommand\labelitemi{{\boldmath$\cdot$}}
\item 
\href{http://soystage.uark.edu}{http://soystage.uark.edu}
\item Webtool that predicts timing of lifecyle
	events for soybeans grown in the
	Midsouthern US.
\end{itemize}



\section*{\sc Professional service and leadership}
\hrule
\vspace{.1in}
{\bf Member of the Curriculum Committee of the
Crop, Soil, and Environmental Sciences Major \hfill  2019}
\vspace{-.1in}
\begin{itemize}
\itemsep-0.5em
\renewcommand\labelitemi{{\boldmath$\cdot$}}
\item 
University of Arkansas,
Fayetteville, Arkansas, US
\end{itemize}

{\bf President of the Crop, Soil, and
Environmental Sciences Graduate Student
Club\hfill  2019} \vspace{-.1in} \begin{itemize}
	\itemsep-0.5em
	\renewcommand\labelitemi{{\boldmath$\cdot$}}
\item University of Arkansas, Fayetteville,
Arkansas, US \end{itemize}

{\bf Vice president of the Crop, Soil, and
Environmental Sciences Graduate Student
Club\hfill  2018}
\vspace{-.1in}
\begin{itemize}
\itemsep-0.5em
\renewcommand\labelitemi{{\boldmath$\cdot$}}
\item University of Arkansas, 
Fayetteville, Arkansas, US
\end{itemize}

\section*{\sc Professional Memberships}
\hrule
\vspace{.1in}
American Society of Agronomy (ASA) \hfill{\bf 2018 - present}

Crop Science Society of America (CSSA) \hfill{\bf 2018 - present}

Soil Science Society of America (SSSA) \hfill{\bf 2018 - present}

International Society of Precision Agriculture \hfill {\bf 2024 - present}

\section*{\sc Languages}
\hrule
\vspace{.1in}
English - Fluent

Portuguese - Native

\section* {\sc Programming languages}
\hrule
\vspace{.1in}
R, Python, C\#, and JavaScript

%% \section* {\sc References}
%% \hrule
%% 
%% \vspace{.1in}
%% {\bf Dr. Fernando Miguez}\\ Professor\\ Iowa State
%% University\\ Department of Agronomy \\ Email:
%% \href{mailto:femiguez@iastate.edu}{femiguez@iastate.edu}
%% \vspace{.1in}
%%
%% {\bf Dr. Trent Roberts}\\ Professor \\ University of
%% Arkansas\\ Department of Crop, Soil, and
%% Environmental Sciences \\ Email:
%% \href{mailto:tlrobert@uark.edu}{tlrobert@uark.edu}
%% 
%% \vspace{.1in} {\bf Dr. Larry Purcell }\\ Emeritus
%% Distinguished Professor\\ University of
%% Arkansas\\ Department of Crop, Soil, and
%% Environmental Sciences \\ Email:
%% \href{mailto:lpurcell@uark.edu}{lpurcell@uark.edu}
 
\end{document}




